\begin{table}[H]
\centering
\footnotesize
\caption{Acierto de género por especie (\%). Se incluyen todas las especies con al menos cinco observaciones y menos de 90\,\% en al menos una evaluación disponible. Un asterisco indica soporte menor de cinco en esa celda; -- indica ausencia de observaciones.}
\label{tab:species}
\begin{tabularx}{\linewidth}{Y r r r r r r r r}
\toprule
\textbf{Especie} & \textbf{Obs. val} & \textbf{Obs. test} & \textbf{R val} & \textbf{A val} & \textbf{B val} & \textbf{C val} & \textbf{D val} & \textbf{R test previo}\\
\midrule
\textit{Amanita franchetii} & 3 & 11 & 100,0$^{*}$ & 100,0$^{*}$ & 100,0$^{*}$ & 100,0$^{*}$ & 100,0$^{*}$ & 81,8\\
\textit{Auricularia mesenterica} & 7 & 13 & 42,9 & 57,1 & 14,3 & 57,1 & 42,9 & 61,5\\
\textit{Ganoderma applanatum} & 177 & 159 & 88,1 & 97,2 & 92,7 & 97,2 & 94,4 & 89,3\\
\textit{Ganoderma lucidum} & 32 & 71 & 93,8 & 96,9 & 96,9 & 96,9 & 96,9 & 88,7\\
\textit{Ganoderma resinaceum} & 14 & 27 & 85,7 & 92,9 & 92,9 & 100,0 & 100,0 & 92,6\\
\textit{Geastrum melanocephalum} & 0 & 7 & -- & -- & -- & -- & -- & 85,7\\
\textit{Geastrum striatum} & 7 & 22 & 100,0 & 85,7 & 100,0 & 100,0 & 85,7 & 95,5\\
\textit{Hericium cirrhatum} & 17 & 15 & 76,5 & 88,2 & 82,4 & 82,4 & 94,1 & 73,3\\
\textit{Lycoperdon mammiforme} & 2 & 6 & 100,0$^{*}$ & 100,0$^{*}$ & 100,0$^{*}$ & 50,0$^{*}$ & 100,0$^{*}$ & 83,3\\
\textit{Lycoperdon molle} & 5 & 6 & 100,0 & 100,0 & 100,0 & 100,0 & 100,0 & 83,3\\
\textit{Morchella esculenta} & 5 & 13 & 100,0 & 100,0 & 100,0 & 100,0 & 100,0 & 84,6\\
\textit{Pleurotus dryinus} & 17 & 19 & 94,1 & 100,0 & 100,0 & 100,0 & 100,0 & 89,5\\
\textit{Trametes cinnabarina} & 26 & 19 & 76,9 & 96,2 & 92,3 & 92,3 & 92,3 & 68,4\\
\bottomrule
\end{tabularx}
\source{reports/s6/comparison/species_difficult.csv}
\end{table}
